% ===== PRÉAMBULE CANONIQUE — à coller VERBATIM en tête de CHAQUE .tex =====
\documentclass[11pt,a4paper]{article}
\usepackage[utf8]{inputenc}\usepackage[T1]{fontenc}\usepackage{lmodern}
\usepackage[french]{babel}
\usepackage{amsmath,amssymb,mathtools}\usepackage{icomma}
\usepackage{siunitx}\sisetup{locale=FR}  % \qty{}{}, \si{}, \ang{}
\usepackage{xcolor}\usepackage{colortbl}
\usepackage{tikz}\usepackage{pgfplots}\pgfplotsset{compat=1.18}
\usepackage[siunitx,europeanresistors]{circuitikz} % circuits (élec.)
\usepackage[most]{tcolorbox}\usepackage{enumitem}\usepackage{array,booktabs}
\usepackage[a4paper,margin=2.2cm]{geometry}
\usetikzlibrary{arrows.meta,calc,angles,quotes,positioning,patterns,%
  backgrounds,shapes.geometric,decorations.pathmorphing,decorations.markings,intersections}
\definecolor{primary}{RGB}{222,49,99}\definecolor{primarydark}{RGB}{150,22,55}
\definecolor{defcol}{RGB}{0,114,178}\definecolor{methcol}{RGB}{52,92,148}
\definecolor{excol}{RGB}{0,135,90}\definecolor{attncol}{RGB}{200,40,55}
\tcbset{boxbase/.style={enhanced,breakable,boxrule=0.9pt,arc=2.5pt,
  left=3mm,right=3mm,top=2.3mm,bottom=2.3mm,fonttitle=\bfseries,coltitle=white}}
% --- boîtes TUTORIEL (noms EXACTS, ne pas renommer) ---
\newtcolorbox{cle}[1][]{boxbase,colback=defcol!6,colframe=defcol,
  title={\textsc{À retenir}\ifx\relax#1\relax\relax\else\textmd{ : \itshape #1}\fi}}
\newtcolorbox{methode}[1][]{boxbase,colback=methcol!7,colframe=methcol,sharp corners,
  title={\textsc{Méthode}\ifx\relax#1\relax\relax\else\textmd{ --- \itshape #1}\fi}}
\newtcolorbox{exemplebox}[1][]{boxbase,colback=excol!5,colframe=excol,
  title={\textsc{Exemple}\ifx\relax#1\relax\relax\else\textmd{ : \itshape #1}\fi}}
\newtcolorbox{piege}[1][]{boxbase,colback=attncol!6,colframe=attncol,
  title={\textsc{Piège}\ifx\relax#1\relax\relax\else\textmd{ : \itshape #1}\fi}}
\newtcolorbox{remarque}[1][]{boxbase,colback=black!4,colframe=black!45,
  title={\textsc{Remarque}\ifx\relax#1\relax\relax\else\textmd{ : \itshape #1}\fi}}
% --- boîtes EXERCICE / CORRIGÉ (noms EXACTS, auto-comptées) ---
\newtcolorbox[auto counter]{exo}[1][]{boxbase,colback=primary!4,colframe=primary,
  title={\textsc{Exercice}~\thetcbcounter\ifx\relax#1\relax\relax\else\textmd{ --- \itshape #1}\fi}}
\newtcolorbox[auto counter]{corrige}[1][]{boxbase,colback=excol!4,colframe=excol,
  title={\textsc{Corrigé}~\thetcbcounter\ifx\relax#1\relax\relax\else\textmd{ --- \itshape #1}\fi}}
\newcommand{\vv}[1]{\vec{#1}}\newcommand{\ds}{\displaystyle}
\newcommand{\R}{\mathbb{R}}\newcommand{\N}{\mathbb{N}}\newcommand{\Z}{\mathbb{Z}}
% =========================================================================
\begin{document}
\sisetup{per-mode=symbol}

\begin{corrige}[la force de réaction]
\begin{enumerate}[label=\alph*)]
  \item Même intensité : $F_{B\to A}=\qty{15}{\newton}$ (3\textsuperscript{e} loi).
  \item Même direction (même droite d'action), \emph{sens opposé}, et elle s'applique sur le corps $A$.
\end{enumerate}
\end{corrige}

\begin{corrige}[deux patineurs se repoussent]
\begin{enumerate}[label=\alph*)]
  \item $a_A=\dfrac{F}{m_A}=\dfrac{60}{60}=\qty{1}{\metre\per\second\squared}$.
  \item $a_B=\dfrac{F}{m_B}=\dfrac{60}{40}=\qty{1,5}{\metre\per\second\squared}$ (le plus léger accélère le plus).
\end{enumerate}
\end{corrige}

\begin{corrige}[rapport des accélérations]
\begin{enumerate}[label=\alph*)]
  \item $\dfrac{a_{\text{boulet}}}{a_{\text{canon}}}=\dfrac{F/m_{\text{boulet}}}{F/m_{\text{canon}}}
        =\dfrac{m_{\text{canon}}}{m_{\text{boulet}}}$.
  \item $=\dfrac{1000}{10}=100$ : le boulet accélère $100$ fois plus que le canon.
\end{enumerate}
\end{corrige}

\begin{corrige}[qui accélère le plus ?]
\begin{enumerate}[label=\alph*)]
  \item La \textbf{petite voiture} : à force $F$ égale, $a=F/m$ est d'autant plus grande que la masse est petite.
  \item Oui, exactement la même intensité (action--réaction) ; seules les accélérations diffèrent.
\end{enumerate}
\end{corrige}

\begin{corrige}[action et réaction dans la vie courante]
\begin{enumerate}[label=\alph*)]
  \item Action : le pied pousse le sol vers l'arrière ; réaction : le sol pousse le marcheur vers l'avant (c'est elle qui fait avancer).
  \item Action : les bras poussent l'eau vers l'arrière ; réaction : l'eau propulse le nageur vers l'avant.
\end{enumerate}
\end{corrige}
\end{document}
